\section{Exact certificates and reproducibility}\label{sec:verification}

The finite row certificates and the scalar inequalities in the analytic
proof are checked using integers and rational numbers. The supplied
verifier uses Python's standard library; it requires no computer algebra
system, numerical root finder, or external package. From the supplement's
root directory, the command is
\begin{verbatim}
python src/verify_all.py
\end{verbatim}
The script also accepts invocation by its absolute path from another
working directory. All component paths are resolved relative to the
script. It checks that the row certificate ends at $250$, checks the
presence of every required verifier, and stops at the first failed
assertion or process. Python must be run without the optimization option
\texttt{-O}, which disables assertions.

\subsection{The finite row certificates}

For $1\le n\le250$, the file
\path{certificates/sign_certificates_n1_250.json} contains $n-1$ rational
numbers
\[
 0>q_{n,1}>q_{n,2}>\cdots>q_{n,n-1}.
\]
The verifier independently computes the Eulerian numbers from
\[
 A(n,k)=\sum_{i=0}^k(-1)^i\binom{n+1}{i}(k-i)^n
       \qquad(1\le k\le n)
\]
and checks agreement with the recurrence. It then forms the matrix-square
coefficients
\[
 B(n,k)=\sum_{j=k}^n A(n,j)A(j,k),\qquad
 R_n(x)=P_n(x)/x.
\]
It checks $B(n,n)=1$, $B(n,1)=n!$, positivity of all coefficients of
$R_n$, the order and negativity of the certificate points, and
\[
 \operatorname{sgn}R_n(q_{n,i})=(-1)^i
       \qquad(1\le i\le n-1).
\]
Since $R_n(0)=n!>0$, these signs give one root in each of the $n-1$
disjoint intervals between consecutive listed points, starting at zero.
They exhaust the degree. Thus every root of $R_n$ is negative and simple;
$P_n$ has the additional simple root at zero. The case $P_0=1$ is immediate.

The sign evaluations clear denominators during Horner evaluation.
For $q=a/d$ with $d>0$ and
$R(x)=\sum_{i=0}^m r_i x^i$, the sign is the sign of the integer
\[
 d^mR(a/d)=\sum_{i=0}^m r_i a^i d^{m-i}.
\]
Starting with $h=r_m$, the verifier successively replaces
$h$ by $ah+r_i d^{m-i}$ for $i=m-1,\ldots,0$.
It retains the original rational Horner routine and checks equality of
the two evaluations, after multiplying by $d^m$, at every certificate
point in the first $20$ rows.
The verifier does not rely on the numerical procedure that originally
selected the rational points.

\subsection{Rational bounds for the analytic constants}

The scalar verifiers use positive Taylor series and geometric remainder
bounds. For rational $x\ge0$ and $M+2>x$, put
\[
 E_M(x)=\sum_{j=0}^{M}\frac{x^j}{j!}.
\]
Then
\[
 E_M(x)\le e^x\le E_M(x)+
 \frac{x^{M+1}}{(M+1)!}\frac1{1-x/(M+2)}.
\]
Taking reciprocals gives outward bounds for $e^{-x}$. To bound logarithms,
write $x=2^r y$ with $1\le y\le2$, and use
\[
 \log y=2\sum_{j=0}^{\infty}\frac{z^{2j+1}}{2j+1},
 \qquad z=\frac{y-1}{y+1}\in[0,1/3].
\]
After terms $0,\ldots,M-1$, the remaining sum is at most
\[
 \frac{2z^{2M+1}}{(2M+1)(1-z^2)}.
\]
The same formula bounds $\log2$, and reciprocation handles $x<1$.
The series
\[
 I_0(v)=\sum_{j=0}^{\infty}\frac{(v^2/4)^j}{(j!)^2}
\]
is treated in the same way, with successive tail ratios bounded by
$(v^2/4)/(M+2)^2$. Square-root comparisons are reduced to squared
rational inequalities with the signs checked first.

For inequalities on unbounded intervals, the proof first reduces the
claim to positivity of a polynomial in a shifted nonnegative variable;
the verifier checks its coefficients. On a bounded interval the
corresponding Bernstein coefficients are used. These are finite rational
calculations. The reduction from the original function to the polynomial
inequality is given in the relevant proof above.

Table~\ref{tab:scalar-summary} records the principal bounds used in the
proof. Every terminating decimal in the table denotes an exact rational
number. The estimates for $n\ge250$ are uniform in $n$; the program checks
the scalar reductions, rather than enumerating values of $n$.

\begin{table}[htbp]
\centering\small
\begin{tabular}{p{.46\textwidth}p{.43\textwidth}}
\toprule
Quantity and range & Certified bound \\
\midrule
Positive pole tail, $0\le s\le3\sqrt n$ & $<.85$ \\
Positive pole tail, $3\sqrt n\le s\le4\sqrt n$ & $<.33$ \\
Negative pole tail, $-4\sqrt n\le s\le0$ & $<.88$ \\
Variance parameters, $\pi\le u\le4\sqrt n$ &
 $Y>7.5$, $Y/n<.26$, $h=Y^2/n<17$ \\
Tilted variance & $\operatorname{Var}L<8/h$ \\
Quadratically weighted far tail in the variance bound & $<.003$ \\
Poisson--binomial maximum mass, $v\ge.7$ & $<.47/\sqrt v$ \\
Each weighted tail outside a paired rectangle & $<.04AB$ \\
Ideal mass outside an isolated mode & $<.907$ times the mode \\
Eulerian replacement error in the overlap band & $<AB/(1000H)$ \\
Independent-zero location error & $<.001$ in $\log X$ \\
Derivative near each independent zero & $>.128$ \\
Paired middle sign after all errors & $<-.79AB/H$ \\
Same-coordinate event spacing in the overlap band & $>12$ \\
\bottomrule
\end{tabular}
\caption{Principal rational constants in the proof.}
\label{tab:scalar-summary}
\end{table}

For example, write
\[
 S(a)=\sum_{j\ge1}j^2e^{-aj(j-1)},\qquad
 F(h)=.47\{S(h/7)+e^{-h/2}S(h/3)+.003\}.
\]
The maximum-mass argument for the variance reduces to
$F(h)<(8/h)^{3/2}$. On each interval $[a,b]$ in
Table~\ref{tab:variance-numerical}, monotonicity bounds $F(h)$ by the
listed rational upper bound for $F(a)$, whose square is less than
$(8/b)^3$. The range $0<h\le2$ is covered by
$S(a)<2a^{-3/2}$ for $a\le2/3$, as proved above.
For $11\le h<17$, the two-adjacent-point argument is checked on the
twelve intervals $[j/2,(j+1)/2]$, $j=22,\ldots,33$. On each interval,
the verifier checks the positive gap and the squared inequality used to
bound $\frac14+\Delta+\sqrt\Delta$ by $8/h$.

\begin{table}[htbp]
\centering\small
\begin{tabular}{ccc}
\toprule
$a$ & $b$ & Upper bound for $F(a)$ \\
\midrule
2&3&2.969\\
3&4&1.803\\
4&6&1.298\\
6&8&.861\\
8&10&.676\\
10&11&.584\\
\bottomrule
\end{tabular}
\caption{The six intermediate intervals in the variance estimate.}
\label{tab:variance-numerical}
\end{table}

For the edge bridge $250\le n\le1600$, put
$L=\lceil\sqrt n/4\rceil$. If $\ell\le L$, then $\ell\le10$ and
\[
 n\ge n_{\min}(\ell)=\max\{250,16(\ell-1)^2+1\}.
\]
The required upper bounds decrease with $n$ at fixed $\ell$, and the
affine lower bounds increase, so the
ten endpoint checks in Table~\ref{tab:finite-edge-endpoints} cover the
whole range. They give outside coordinate ratios
$(.036,.015,.0006,.015)$ in the four directions, total ideal error
$<.054$, and central Eulerian error $<10^{-6}$. Their sum is below the
quadratic sign margin $1/15$. The same checks give the phase comparison
and $40<\log X_L^-<4\sqrt n$.

\begin{table}[htbp]
\centering\small
\begin{tabular}{cccccccc}
\toprule
$\ell$&$1,2,3,4$&5&6&7&8&9&10\\
\midrule
$n_{\min}(\ell)$&250&257&401&577&785&1025&1297\\
\bottomrule
\end{tabular}
\caption{Endpoints used for the finite edge bridge.}
\label{tab:finite-edge-endpoints}
\end{table}

\subsection{Supplement map}

The reproduction entry point is \path{src/verify_all.py}. The required
files are organized as follows.
\begin{itemize}
\item \path{src/verify_certificates.py} and
\path{src/eulerian_core.py} reconstruct the Eulerian rows and check the
rational sign certificates.
\item Under \path{src/astra_literature/},
\path{verify_extended_variance_constants.py} checks the variance estimate,
\path{verify_joint_tail_audit_constants.py} checks both lattice tails and
the Eulerian replacement error, and \path{verify_finite_edge_bridge.py}
checks the finite-range edge and phase constants.
\item Under \path{src/astra_structural/},
\path{verify_positive_bulk_constants.py} and
\path{verify_positive_bulk_extension_constants.py} check the two positive
bulk ranges, while \path{verify_negative_bulk_pole_constants_4sqrt.py}
checks the negative bulk through $4\sqrt n$.
These scripts import their rational Taylor helpers from
\path{verify_negative_bulk_variance_constants.py}, which is also included.
The file \path{verify_paired_cluster_sign_constants.py} checks the
independent-zero and paired-sign constants.
\item Under \path{src/astra_interlacing/},
\path{verify_mode_geometry_constants.py} checks the overlap-band
geometry, outer separator, lattice-sum margins, and separator phase
constants.
\end{itemize}

The finite polynomial certificates and the scalar checks have distinct
roles. The former prove the initial rows directly. The latter verify the
explicit constants in the uniform analytic proof, whose root count
covers all remaining rows.
